%author : Jiahao Shi 2024/5/10
\documentclass[10pt]{beamer}
\usepackage[utf8]{inputenc}
\usepackage{xeCJK}
\usepackage{graphicx}
\usepackage {mathtools}
\usepackage{utopia} %font utopia imported
\usetheme{CambridgeUS}
\usecolortheme{dolphin}
\usepackage{tcolorbox}
\usepackage{tabularray}
\usepackage{float}
\usepackage{adjustbox}
\usepackage{caption}
\usepackage{subcaption}
\usepackage{multirow}
\usepackage{tikz}
\usepackage{booktabs}
\usepackage{makecell}
\usepackage{array}
\renewcommand\theadfont{\bfseries}

\usetikzlibrary{positioning} % For advanced node positioning (right=of)
\usetikzlibrary{calc}        % For coordinate calculations
\usetikzlibrary{arrows.meta}  % For improved arrow tips
\usetikzlibrary{shapes} 
\usepackage{longtable}
% custom colors 自定义颜色，可根据主题调整
% Define a palette of reds
% Define redgrades first (example)
\definecolor{lightred}{RGB}{180,227,236}        % very light red (backgrounds)
\definecolor{middlelightred}{RGB}{138,22,57}  % light/mid red (accents)
\definecolor{middledered}{RGB}{217,79,120}      % bright medium red (main highlight)
\definecolor{redforsnnu}{RGB}{176,44,89}        % deep but still clear red
\definecolor{darkred}{RGB}{138,22,57}           % uni logo red (titles, emphasis)


% Assign redgrades to Beamer elements
\setbeamercolor*{palette primary}{bg=redforsnnu, fg=white}       % main deep red background
\setbeamercolor*{palette secondary}{bg=middledered, fg=white}    % medium red background
\setbeamercolor*{palette tertiary}{bg=middlelightred, fg=white} % lighter red background, dark text
\setbeamercolor*{titlelike}{fg=redforsnnu}                       % deep red foreground for title-like elements
\setbeamercolor*{title}{bg=darkred, fg=white}                    % very dark red background with white text
\setbeamercolor*{item}{fg=redforsnnu}                            % deep red for items
\setbeamercolor*{caption name}{fg=middledered}                   % medium red for caption names


\usefonttheme{professionalfonts}
\usepackage{natbib}
\usepackage{hyperref}

\usepackage{tikz}
\usepackage{graphicx}
\usepackage{lipsum}
\usetikzlibrary{arrows.meta, shapes.geometric}

\titlegraphic{
\begin{tikzpicture}[overlay, remember picture]
    % Position the left logo
    %\node[anchor=north west, xshift=0.2cm, yshift=-0.2cm] at (current page.north west){
        
    %};
    % Add text in between the logos
    \node[anchor=north, yshift=-0.7cm, xshift =-0cm] at (current page.north){
        \begin{minipage}{\textwidth}
            \centering
             \includegraphics[height=2cm]{UniversiteParisCite_logo_horizontal_couleur_RVB.jpg}
        \end{minipage}
    };
    % Position the right logo
    %\node[anchor=north east, xshift=-0.2cm, yshift=-0.2cm] at (current page.north east){
    %};
\end{tikzpicture}
}

\setbeamerfont{title}{size=\large}
\setbeamerfont{subtitle}{size=\small}
\setbeamerfont{author}{size=\small}
\setbeamerfont{date}{size=\small}
\setbeamerfont{institute}{size=\small}

\vspace*{0.5cm}
\title[]{Extraction de connaissances à partir de textes non structurés à l'aide de techniques d'apprentissage automatique}

% \subtitle{ } % Subtitle (if needed)
\author[\textbf{Recherche et extraction sémantique à partir de texte}]{
            { \bf BOUKHARI Imed Eddine}\\ 
            {\bf OMARI Hamza} \\ 
            {\bf AGOUNI KACI Sofiane} \\ 
            {\bf SOUAYAH Abdelkader} } % Author
\vspace*{-0.5cm}
\date[\textcolor{white}{18/12/2025}]{18/12/2025}

% Insert the middle text
\vspace*{2cm}

\vspace*{-1cm}

%------------------------------------------------------------
%This block of commands puts the table of contents at the 
%beginning of each section and highlights the current section:

\AtBeginSection[]{
  \begin{frame}
  \vfill
  \transfade
  \centering
  \begin{beamercolorbox}[sep=8pt,center,shadow=true,rounded=true]{title}
    \usebeamerfont{title}\insertsectionhead\par%
  \end{beamercolorbox}
  \vfill
  \end{frame}
}
%------------------------------------------------------------
\begin{document}

\frame{\titlepage}

\begin{frame}{Table of Contents}
    \tableofcontents
\end{frame}

\section{Dataset}

%------------------------------------------------------------
% SLIDE 1: The Context & Problem
%------------------------------------------------------------
\begin{frame}{Ensemble de données}
        \vspace{-0.6cm}
    \begin{columns}[T]
        % RIGHT COLUMN - TEXT
        \begin{column}{0.38\textwidth}
            \begin{tcolorbox}[colback=lightred!20, colframe=redforsnnu, title=Aperçu de l'ensemble de données]
                \begin{itemize}\scriptsize
                    \item 17 classes
                    \item $\sim$ 1 000 000 mots
                    \item $\sim$ 500 000 phrases
                \end{itemize}
            \end{tcolorbox}

            
            \begin{tcolorbox}[colback=middledered!20, colframe=darkred, title=Étiquettes d'entités]
                \begin{itemize}\scriptsize
                    \item \textbf{O}: Aucune étiquette (87\%)
                    \item \textbf{geo}: Entité géographique
                    \item \textbf{gpe}: Entité géopolitique
                    \item \textbf{per}: Personne
                    \item \textbf{org}: Organisation
                    \item \textbf{tim}: Expression temporelle
                    \item \textbf{art}: Artefact
                    \item \textbf{nat}: Nationalité
                    \item \textbf{eve}: Événement
                \end{itemize}
            \end{tcolorbox}
        \end{column}
        
% LEFT COLUMN - IMAGE
        \begin{column}{0.58\textwidth}
            \begin{figure}
                \centering
                \includegraphics[width=\linewidth]{NER-Tags.png}
            \end{figure}
            \begin{tcolorbox}[colback=lightred!20, colframe=redforsnnu, title=]
                \begin{itemize}\scriptsize                    \item Événements géopolitiques (Guerre en Irak, ..)
                    \item Attributs d'entités: Personne, Org, ..
                \end{itemize}
            \end{tcolorbox}
        \end{column}
        
    \end{columns}

    
\end{frame}




\section{Preprocessing}
\begin{frame}{NLP preprocessing pipelines }
    
    \begin{figure}
        \centering
        \includegraphics[width=0.6\linewidth]{generic.png}
        \caption*{Generic NLP Pipeline}
        \label{fig:generic}
    \end{figure}
    
    \vspace{-0.5cm}
    \begin{figure}
        \centering
        \includegraphics[width=0.6\linewidth]{DL.png}

    \end{figure}
    
\end{frame}

\begin{frame}{Preprocessing}    
    \begin{columns}[T]
        % LEFT COLUMN - Data Cleaning
        \begin{column}{0.48\textwidth}
            \begin{figure}
                \centering
                \includegraphics[width=\linewidth]{datacleaning.png}
                \caption*{Data Cleaning}
                \label{fig:datacleaning}
            \end{figure}
        \end{column}
        
        % RIGHT COLUMN - Best Free NLP Tools
        \begin{column}{0.48\textwidth}
            \begin{figure}
                \centering
                \includegraphics[width=\linewidth]{bestfreenlptools.png}
                \caption*{Best Free NLP Tools}
                \label{fig:besttools}
            \end{figure}
        \end{column}
    \end{columns}
    
\end{frame}
\begin{frame}{Text Representation}
        
    \begin{columns}[T]
        % LEFT COLUMN - Hand crafted features
        \begin{column}{0.48\textwidth}
        \vspace{0.5cm}
            \textbf{Hand crafted features}
            \begin{itemize}\small
                \item Tokenisation
                \item Balisage grammatical (POS tagging)
                \item Lemmatisation
                \item Reconnaissance d'entités nommées (NER)
            \end{itemize}

        \end{column}
        
        % RIGHT COLUMN - Images
        \begin{column}{0.48\textwidth}
            \vspace{-0.5cm}
            \begin{figure}
                \centering
                \includegraphics[width=\linewidth]{word2vec.png}
                \caption*{Word2Vec}
                \label{fig:word2vec}
            \end{figure}
            
            \vspace{-0.5cm}
            
            \begin{figure}
                \centering
                \includegraphics[width=\linewidth]{bert.png}
                \caption*{BERT}
                \label{fig:bert}
            \end{figure}
        \end{column}
    \end{columns}
    
\end{frame}

\begin{frame}{Text Representation -2-}
    \begin{figure}
        \centering
        \includegraphics[width=1\linewidth]{fig.png}
    \end{figure}
\end{frame}

\section{Classification NER}
\begin{frame}{Méthodes Implémentées}
    
    \begin{columns}[T]
        % LEFT COLUMN - Rule-based
        \begin{column}{0.48\textwidth}
            \begin{tcolorbox}[colback=lightred!20, colframe=redforsnnu, title=Rule-based NER]
                \begin{itemize}
                    \item \textbf{Utilisation:} spaCy
                    \item Règles linguistiques + modèles pré-entraînés
                    \item \textbf{Avantages:}
                    \begin{itemize}
                        \item Rapide à mettre en place
                        \item Interprétable
                    \end{itemize}
                    \item \textbf{Limites:}
                    \begin{itemize}
                        \item Sensible au bruit
                        \item Faible généralisation
                    \end{itemize}
                \end{itemize}
            \end{tcolorbox}
        \end{column}
        
        % RIGHT COLUMN - Machine Learning
        \begin{column}{0.48\textwidth}
            \begin{tcolorbox}[colback=middledered!20, colframe=darkred, title=Machine Learning NER]
                \begin{itemize}
                    \item \textbf{Modèle:} CRF
                    \item Apprentissage supervisé à partir du corpus
                    \item \textbf{Avantages:}
                    \begin{itemize}
                        \item Meilleure précision
                        \item Prise en compte du contexte
                    \end{itemize}
                    \item \textbf{Limites:}
                    \begin{itemize}
                        \item Dépendance à la qualité des données annotées
                    \end{itemize}
                \end{itemize}
            \end{tcolorbox}
        \end{column}
    \end{columns}
\end{frame}

 
\begin{frame}{Objectif \& Données}
    \vspace{-0.5cm}
    
    \begin{tcolorbox}[colback=lightred!20, colframe=redforsnnu, title=]
        \begin{itemize}
            \item \textbf{Objectif:} Appliquer et comparer plusieurs approches de Named Entity Recognition (NER)
            \item Identifier et classifier des entités nommées dans un corpus annoté
            \item \textbf{Données:} Corpus NER annoté (schéma BIO)
            \item Colonnes: Word, POS, NER Tag
            \item Catégories d'entités: Person, Organization, Location, Date, Product, etc.
            \item \textbf{But final:} Extraire des entités fiables pour des tâches d'extraction d'information et de relations
        \end{itemize}
    \end{tcolorbox}
    
\end{frame}

\begin{frame}{Résultats \& Évaluation}

    
    \begin{columns}[T]
        % LEFT COLUMN - TEXT
        \begin{column}{0.33\textwidth}
            \begin{itemize}\small
                \item \textbf{Classification des entités}
                \begin{itemize}\small
                    \item Chaque entité est classée dans une catégorie (Person, Organization, Date, etc.)
                    \item Utilisation du schéma BIO pour la prédiction séquentielle
                \end{itemize}
                \item \textbf{Évaluation}
                \begin{itemize}\small
                    \item Métriques: Precision: \textbf{0.86},\\ Recall: \textbf{0.84}, \\ F1-score: \textbf{0.85}.
                \end{itemize}
            \end{itemize}
        \end{column}
        
        % RIGHT COLUMN - IMAGE
        \begin{column}{0.65\textwidth}
            \begin{figure}
                \centering
                \includegraphics[width=\linewidth]{hmz.png}
            \end{figure}
        \end{column}
    \end{columns}
    
\end{frame}

\begin{frame}{Observations}
    
    \begin{itemize}
        \item Les approches ML surpassent les méthodes rule-based
        \item Le CRF offre un bon compromis performance / complexité
        \item Les erreurs proviennent principalement:
        \begin{itemize}
            \item du déséquilibre des classes
            \item des ambiguïtés linguistiques
        \end{itemize}
    \end{itemize}
    
\end{frame}

\section{Extraction des relations et le Knowledge graph}

% SLIDE 1 – Extraction des relations
\begin{frame}{Extraction des relations}
    \vspace{-0.4cm}
    \begin{itemize}
        \item \textbf{Objectif:} extraire des triplets \textit{(sujet, verbe, objet)} à partir du corpus.
        
        \item \textbf{Étape 1 – Sélection des verbes}
        \begin{itemize}
            \item On ne garde que des verbes dans une liste \texttt{top\_verbs}
            \item Exemple: \textit{say, meet, visit, arrest, attack, call, report, warn, ...}
        \end{itemize}
        
        \item \textbf{Étape 2 – Parcours des dépendances}
        \begin{itemize}
            \item \texttt{nsubj}, \texttt{nsubjpass} pour trouver le sujet du verbe  
                  Exemple: \textit{"Clinton visited Iraq"} $\rightarrow$ \texttt{nsubj(visited, Clinton)}
            \item \texttt{dobj}, \texttt{pobj}, pour trouve l’objet  
                  Exemple: \textit{"visited Iraq"} $\rightarrow$ \texttt{dobj(visited, Iraq)}
        \end{itemize}
        
        \item \textbf{Étape 3 – Construction des triplets}
        \begin{itemize}
            \item Pour chaque verbe sélectionné, on génère \textit{(sujet, verbe, objet)}
            \item Exemple de triplet: \textbf{(Clinton, visit, Iraq)}
        \end{itemize}
        
        \item \textbf{Étape 4 – Chunking du texte}
        \begin{itemize}
            \item Texte très long $\rightarrow$ découpage en blocs de 75 000 caractères
            \item Extraction des triplets sur chaque bloc, puis agrégation des résultats
        \end{itemize}
        
    \end{itemize}
\end{frame}


% SLIDE 2 – Knowledge graph sans coréférence

\begin{frame}{Knowledge graph (sans coréférence)}
    \vspace{-0.4cm}
    \begin{itemize}\small
        \item Les triplets \textit{(sujet, verbe, objet)} sont transformés en graphe orienté avec NetworkX.
        \item Nœuds: entités nommées; arêtes: relations verbales (étiquette = verbe).
        \item Visualisation: \texttt{spring\_layout}, nœuds bleus, labels d'arêtes en rouge.
    \end{itemize}
    
    \begin{figure}
        \centering
        \includegraphics[width=0.75\linewidth]{kg\_no\_coref.png} % <-- mets ici le nom de la 1ère image
        \caption*{Knowledge Graph – relations extraites (avant coréférence)}
    \end{figure}
\end{frame}

% SLIDE 3 – Coréférence et nouveau graphe

\begin{frame}{Coréférence et Knowledge graph}
    \vspace{-0.4cm}
    \begin{itemize}\small
        \item Coréférence optimisée pour longs documents ($30\,000+$ mots).
        \item Traitement phrase par phrase avec fenêtre de contexte glissante.
        \item Remplacement des pronoms (\textit{he, she, it, they, ...}) par:
        \begin{itemize}\scriptsize
            \item entités PERSON, sinon ORG, sinon GPE/LOC.
        \end{itemize}
        \item Les relations sont ré-extraites sur le texte résolu, puis le graphe est reconstruit.
    \end{itemize}
    
    \begin{figure}
        \centering
        \includegraphics[width=0.72\linewidth]{kg\_coref.png} % <-- mets ici le nom de la 2ème image
        \caption*{Knowledge Graph – relations extraites (avec coréférence)}
    \end{figure}
\end{frame}





\begin{frame}{Défis de l'extraction de relations - Analyse et conclusion}
            \vspace{-0.1cm}

    % TABLE AT TOP - FULL WIDTH
    \begin{table}\scriptsize
        \centering
        \begin{tabular}{|p{3cm}|p{3cm}|p{4cm}|}
            \hline
            \textbf{Critère} & \textbf{Analyse de dépendances} & \textbf{RE basée sur l'AM} \\
            \hline
            Approche & Règles \& structures syntaxiques & Apprentissage à partir des données \\
            \hline
            Données annotées requises & Non & Oui \\
            \hline
            Interprétabilité & Élevée & Moyenne à faible \\
            \hline
            Robustesse & Faible & Élevée \\
            \hline
            Généralisation & Limitée & Forte \\
            \hline
            Relations implicites & Non & Oui \\
            \hline
            Scalabilité & Moyenne & Élevée \\
            \hline
            Adaptation à nouveaux domaines & Règles manuelles & Réentraînement automatique \\
            \hline
            Coût computationnel & Faible à modéré & Modéré à élevé \\
            \hline
            Gestion des paraphrases & Faible & Bonne \\
            \hline
            Utilisation typique & Domaines contrôlés & Grands volumes de données, langage varié \\
            \hline
        \end{tabular}
    \end{table}
    
        \vspace{-0.1cm}
    % THREE BOXES BELOW IN COLUMNS
    \begin{columns}[T]
        % BOX 1 - Ambiguity
        \begin{column}{0.31\textwidth}
            \begin{tcolorbox}[colback=lightred!20, colframe=redforsnnu, title=Ambiguïté contextuelle]
                \begin{itemize}\scriptsize
                    \item « Paris » $\rightarrow$ Ville? Personne? Hôtel?
                \end{itemize}
            \end{tcolorbox}
        \end{column}
        
        % BOX 2 - Multi-word entities & Coreference
        \begin{column}{0.31\textwidth}
    \begin{tcolorbox}[colback=middledered!20, colframe=darkred, title=Multi-mots]
        \begin{itemize}\scriptsize
            \item[$\times$] {\color{red}[United] [Nations] [Security] [Council]}
            \item[$\checkmark$] {\color{black}[United Nations Security Council]}
        \end{itemize}
    \end{tcolorbox}
\end{column}
        
        % BOX 3 - Applications & Results
        \begin{column}{0.31\textwidth}
            \begin{tcolorbox}[colback=lightred!20, colframe=redforsnnu, title=Applications]
                \begin{itemize}\scriptsize
                    \item Questions-réponses
                    \item Recommandation
                    \item Recherche sémantique
                    \item \textbf{LLM: 92 vs 40}
                \end{itemize}
            \end{tcolorbox}
        \end{column}
    \end{columns}
    
\end{frame}


\section{Références}

\begin{frame}[allowframebreaks]{Références}
    \begin{thebibliography}{10}
    
        \bibitem{conll2003}
        Sang, E. F. T. K., \& De Meulder, F. (2003).
        ``Introduction to the CoNLL-2003 Shared Task: Language-Independent Named Entity Recognition.''
        \textit{Proceedings of the Seventh Conference on Natural Language Learning}.
        
        \bibitem{spacy}
        Honnibal, M., \& Montani, I. (2017).
        ``spaCy 2: Natural language understanding with Bloom embeddings, convolutional neural networks and incremental parsing.''
        Retrieved from \url{https://spacy.io}
        
        \bibitem{networkx}
        Hagberg, A. A., Schult, D. A., \& Swart, P. J. (2008).
        ``Exploring network structure, dynamics, and function using NetworkX.''
        \textit{Proceedings of the 7th Python in Science Conference}, 11-15.
        
        \bibitem{dataset}
        ``NER Dataset: Named Entity Recognition Corpus.''
        Available on Kaggle and GitHub repositories.
        Corpus with 1,000,000+ words, 500,000+ phrases, 17 entity classes (BIO schema).
        
    \end{thebibliography}
\end{frame}

\end{document}